\section{Reactive-Area Construction and Capacity Mapping}

In three dimensions, the reconstruction extracts an isosurface and deposits
triangle areas into their owning cells. In quasi-two-dimensional images, it
extracts an isoline in the image plane, multiplies its length by the physical
out-of-plane thickness, and assigns that area to the corresponding DBS cells.
The two routes therefore return the same field type and units,
m$^2$\,m$^{-3}_{\mathrm{bulk}}$. The field is constructed before flow and is
held fixed throughout reactive transport.

Intrinsic densities inferred per phase volume or grain mass must first be
converted to a bulk-volume basis, for example using phase fraction and grain
density. Porosity is not a conversion factor for a field already expressed per
bulk volume.

PhreeqcRM \citep{parkhurst2015phreeqcrm} receives area per representative bulk
litre. Accordingly,
$a_{\mathrm{surf}}$ in m$^2$\,m$^{-3}$ is divided by $10^3$ and is not divided
by the mobile-water fraction. For surface master $k$, multiplication by the
site density $N_k$ (mol\,m$^{-2}$) gives the site capacity per representative
bulk volume. OpenFOAM stores amount in kmol, so a dictionary site-density field
in kmol\,m$^{-2}$ is multiplied numerically by $10^3$ when interpreted as
mol\,m$^{-2}$. The units then compose as
\[
\left[\frac{\mathrm{m}^2}{\mathrm{m}^3_{\mathrm{bulk}}}\right]
\left[\frac{\mathrm{mol}}{\mathrm{m}^2}\right]
=\left[\frac{\mathrm{mol}}{\mathrm{m}^3_{\mathrm{bulk}}}\right].
\]

The chemistry mask threshold is set by \texttt{epsChemistryMin} in
\texttt{constant/thermoPhysicalProperties}, and every calculation reported here
uses $\theta_{\mathrm{mask}} = 10^{-3}$.

\section{Discrete PhreeqcRM Coupling}

The surface-complexation cases select \texttt{mixtureType phreeqcMixture} in
\texttt{constant/thermoPhysicalProperties}. The active mapping between
finite-volume cells and PhreeqcRM cells is established once and does not change
during the run. Reconstruction is needed only when a threshold crossing changes
the active-cell map or its topology, or when the fixed area capacity changes.

Each OpenFOAM process constructs PhreeqcRM with one worker thread. Domain
decomposition provides process-level parallelism, and \texttt{runCells()}
advances only the chemistry cells held by that rank. Each case is decomposed
according to the \texttt{NPX}, \texttt{NPY}, and \texttt{NPZ} entries in its
\texttt{system} directory. Every case uses its case-local
\texttt{constant/GeoChem.dat} database. The image-based cases use kinematic
viscosity $\nu = 10^{-6}$~m$^2$\,s$^{-1}$ and pressure drop
$\Delta p = 5\times10^{-5}$~Pa, and write output every 10~s for the 1D
benchmark, 2000~s for the fracture case, and 20~s for the Estaillades case.

\subsection{Checkpoint and Restart Lifecycle}

At an output time, a checkpoint-enabled run first invalidates the old
completion marker collectively, writes a new rank-local dump, writes the
finite-volume fields, and publishes a completion marker only after the field
write succeeds. A crash before publication therefore leaves the partially
updated fields and chemistry state without a marker, and restoration rejects
them. The PhreeqcRM dump and marker must be retained with the corresponding
finite-volume time directory; the \texttt{Surf\_*} fields alone do not
reconstruct the complete internal state.

\section{Short Two-Scale Reconstruction Sensitivity}

The two-scale calculation is a matched short-interval sensitivity test between
two DBS cases that differ only in the surface field. Both use a
$3600\times1200\times1$ grid with $5~\mu$m cells,
16 MPI ranks, $Pe=1$, $D=10^{-9}$~m$^2$\,s$^{-1}$, and prescribed Darcy velocity
$5.46297\times10^{-5}$~m\,s$^{-1}$. The surface field is the only input that
differs, and the other 495 input files were verified as identical.

Both short cases use their case-local \texttt{GeoChem.dat} together with the
same illustrative Li-selective sorbent chemistry. The surface reactions are
$\mathrm{Sor\_aH}+\mathrm{Li^+}\rightleftharpoons\mathrm{Sor\_aLi}+\mathrm{H^+}$
($\log K=-5.8$), and the corresponding Na$^+$, K$^+$, and Mg$^{2+}$ reactions
have $\log K=-7.2$, $-7.0$, and $-7.4$, respectively. Initially Li-free brine
contains 0.10~M Na, 0.01~M K, 0.01~M Mg, $10^{-12}$~M Li, and 0.13~M Cl at pH
7; the inflow changes Li to $10^{-3}$~M and Cl to 0.131~M. These parameters
apply to the short $Pe=1$ sensitivity test only.

Three common checkpoints at 0.0058517566, 0.11703513, and 0.99564208~s meet
the predefined material-response limits. The maximum absolute differences
are $0.000647\%$ in integrated area, $0.000703\%$ in aqueous Li inventory,
$0.2110\%$ in sorbed Li inventory, $0.00172\%$ in aqueous normalised L1,
$0.5147\%$ in sorbed normalised L1, and $0.03993\%$ in fixed-reference-area
sorbed-response L2. Aqueous and sorbed $x_{50}$ shifts are zero. The native-cell
sorbed inventory-density L2 is $3.61$ to $4.49\%$, dominated by the deliberate
relocation of area, so the agreement is in integrated response rather than
pointwise. The last accepted checkpoint covers $0.489\%$ of the configured
$203.53$~s horizon, so the test establishes early material-response agreement.

\section{Reproducibility and Provenance Entry Points}

The paper directory contains self-contained case definitions and source-data
references. From the repository root, the provenance verifier is run with
\begin{verbatim}
papers/scalarTransportDBSFoam-phreeqc-surface-complexation/
    validate_accepted_nonkelvin_artifacts.sh
\end{verbatim}
This command runs the verifier tests and validates
\texttt{figures/accepted\_nonkelvin\_provenance.json}. The manifest binds the
reported 1D, fracture, and Estaillades figures to their source data, renderers,
sidecars, and evidence records by path, size, and SHA-256.

The case-level entry point for the 1D benchmark is
\texttt{runPhase8AnalyticalBenchmark.sh} in
\nolinkurl{cases/1d_linear_retardation}. The fracture and Estaillades cases are
run through the \texttt{runAll.sh} scripts in
\nolinkurl{cases/2d_4StraightFractures} and \nolinkurl{cases/3d_estaillades},
each of which also provides separate mesh, flow, reactive-transport,
processing, and cleanup scripts. The short
two-scale record is in the \texttt{results} directory of
\nolinkurl{cases/two_scale_micromodel_sorption_validation}; it records exact
checkpoints, thresholds, collection hashes, runtime identity, and analyser
identity. The paper sources and both PDFs are built with \texttt{build.sh}. The
flat Springer submission archive is generated transactionally with
\texttt{export\_submission.sh OUTPUT.zip}; its manifest and SHA-256 inventory
are included in the archive.
